\section{Introduction}
\label{sec:intro}

Financial asset recommendation (FAR) sits at the intersection of personalised ranking and regulatory obligation. The European Markets in Financial Instruments Directive (MiFID II) requires firms to assess suitability before recommending a product: alignment between the product's risk profile and the customer's declared risk tolerance, investment horizon, and loss capacity. A recommender that maximises predicted return without conditioning on the customer's risk band can be an excellent ranker on standard benchmarks but a regulatory liability.

FAR-Trans~\cite{sanzcruzado:fartrans} is the first open FAR benchmark with both customer profile metadata and a multi-year transaction history: 228{,}913 Buys across 806 assets and 29{,}090 retail customers, with MiFID II risk bands per customer and metadata sufficient to derive each asset's risk class. Of these, 228{,}241 Buys are \emph{scoreable}: both the customer and the asset have an assigned risk band, enabling pairwise discordance measurement. The accompanying paper benchmarks LightGCN~(a graph-based collaborative filter that propagates embeddings over the user-item bipartite graph), Random Forest, and others on normalised Discounted Cumulative Gain (nDCG) and Return on Investment (ROI), each evaluated on the top-10 ranked items. We additionally report Recall alongside these as a ranking sanity check. None of those metrics measures profile alignment, and none of the published baselines conditions on the declared risk band.

This paper introduces a profile-alignment lens for FAR and applies it in four steps: \emph{diagnose} the prevalence of discordance, \emph{test} whether it carries an economic penalty, \emph{audit} where existing baselines sit on the new axis, and \emph{remedy} the gap with a profile-coherent training objective. These map onto four research questions:

\begin{itemize}
    \item \textbf{RQ1 (Diagnose)}: How prevalent is profile-discordance in observed FAR-Trans Buy transactions, and is it a stable customer trait or transaction-level noise?
    \item \textbf{RQ2 (Test the economic objection)}: Do profile-discordant Buys earn lower realised six-month return than coherent ones, once asset volatility, customer segment, and year are controlled for?
    \item \textbf{RQ3 (Audit existing baselines)}: At their best published settings, do the chosen FAR-Trans baseline models lift profile coherence above the band-conditional random baseline $\pi(b)$ (the expected coherence of a uniformly random recommender) for every declared band?
    \item \textbf{RQ4 (Remedy)}: Can a stratified, profile-coherent LightGCN improve profile coherence over the global baseline, and at what cost to nDCG, Recall, and ROI?
\end{itemize}

\paragraph{Contributions.} This paper makes two contributions. First, we introduce Profile Coherence at top-$k$ ranked items (PC@$k$), a coherence metric with a band-conditional random baseline $\pi(b)$ and scale-invariant lift normalisation. Second, we propose a stratified, profile-coherent LightGCN with a margin loss atop Bayesian Personalised Ranking (BPR, the standard pairwise loss for LightGCN), together with an ablation that isolates the coherence-loss effect from the stratified architecture.
